\hypertarget{ej4}{}
\noindent \textbf{Ejercicio 4.} Determinar si las siguientes fórmulas son tautologías, contradicciones o contingencias.

\begin{multicols}{2}
\begin{enumerate}[label=\alph*)]
    \item $(p \vee \neg p)$
    \item $(p \wedge \neg p)$
    \item $((\neg p \vee q) \leftrightarrow (p \rightarrow q))$
    \item $((p \wedge q) \rightarrow p)$
    \item $((p \wedge (q \vee r)) \leftrightarrow ((p \wedge q) \vee (p \wedge r)))$
    \item $((p \rightarrow (q \rightarrow r)) \rightarrow ((p \rightarrow q) \rightarrow (p \rightarrow r)))$
\end{enumerate}
\end{multicols}

\vspace{0.2cm}
{\color{gray}\hrule}
\vspace{0.4cm}

\textbf{Resolución:}

\begin{enumerate}[label=\textbf{\alph*)}]

    \item \textbf{$(p \vee \neg p)$} \\
    Armamos la tabla de verdad clásica:
    \begin{center}
    \begin{tabular}{|c|c|c|}
    \hline
    $p$ & $\neg p$ & $(p \vee \neg p)$ \\
    \hline
    V & F & \textbf{V} \\
    F & V & \textbf{V} \\
    \hline
    \end{tabular}
    \end{center}
    Como todos los resultados son verdaderos independientemente del valor de $p$, es una \underline{\textbf{Tautología}}.

    \vspace{0.4cm}
    \item \textbf{$(p \wedge \neg p)$}
    \begin{center}
    \begin{tabular}{|c|c|c|}
    \hline
    $p$ & $\neg p$ & $(p \wedge \neg p)$ \\
    \hline
    V & F & \textbf{F} \\
    F & V & \textbf{F} \\
    \hline
    \end{tabular}
    \end{center}
    Como todos los resultados son falsos, es una \underline{\textbf{Contradicción}}.

    \vspace{0.4cm}
    \item \textbf{$((\neg p \vee q) \leftrightarrow (p \rightarrow q))$}
    \begin{center}
    \begin{tabular}{|c|c|c|c|c|}
    \hline
    $p$ & $q$ & $(\neg p \vee q)$ & $(p \rightarrow q)$ & Equivalencia ($\leftrightarrow$) \\
    \hline
    V & V & V & V & \textbf{V} \\
    V & F & F & F & \textbf{V} \\
    F & V & V & V & \textbf{V} \\
    F & F & V & V & \textbf{V} \\
    \hline
    \end{tabular}
    \end{center}
    Ambos lados tienen los mismos valores de verdad, por lo tanto la equivalencia es una \underline{\textbf{Tautología}}.

    \vspace{0.4cm}
    \item \textbf{$((p \wedge q) \rightarrow p)$}
    \begin{center}
    \begin{tabular}{|c|c|c|c|}
    \hline
    $p$ & $q$ & $(p \wedge q)$ & $((p \wedge q) \rightarrow p)$ \\
    \hline
    V & V & V & \textbf{V} \\
    V & F & F & \textbf{V} \\
    F & V & F & \textbf{V} \\
    F & F & F & \textbf{V} \\
    \hline
    \end{tabular}
    \end{center}
    Como no hay ningún caso donde el antecedente sea verdadero y el consecuente falso, la implicación siempre se cumple. Es una \underline{\textbf{Tautología}}.

    \vspace{0.4cm}
    \item \textbf{$((p \wedge (q \vee r)) \leftrightarrow ((p \wedge q) \vee (p \wedge r)))$} \\
    Esta es la propiedad distributiva de la conjunción sobre la disyunción. Evaluando las 8 combinaciones posibles para $p, q, r$, vemos que el bicondicional siempre se cumple:
    \begin{center}
    \begin{tabular}{|c|c|c|c|c|c|}
    \hline
    $p$ & $q$ & $r$ & Lado Izq. & Lado Der. & ($\leftrightarrow$) \\
    \hline
    V & V & V & V & V & \textbf{V} \\
    V & V & F & V & V & \textbf{V} \\
    V & F & V & V & V & \textbf{V} \\
    V & F & F & F & F & \textbf{V} \\
    F & V & V & F & F & \textbf{V} \\
    F & V & F & F & F & \textbf{V} \\
    F & F & V & F & F & \textbf{V} \\
    F & F & F & F & F & \textbf{V} \\
    \hline
    \end{tabular}
    \end{center}
    Es una \underline{\textbf{Tautología}}.

    \vspace{0.4cm}
    \item \textbf{$((p \rightarrow (q \rightarrow r)) \rightarrow ((p \rightarrow q) \rightarrow (p \rightarrow r)))$} \\
    Nuevamente, al evaluar las 8 combinaciones posibles, el resultado de la implicación principal nunca tiene un antecedente $V$ con consecuente $F$.
    \begin{center}
    \begin{tabular}{|c|c|c|c|c|c|}
    \hline
    $p$ & $q$ & $r$ & Lado Izq. & Lado Der. & ($\rightarrow$) \\
    \hline
    V & V & V & V & V & \textbf{V} \\
    V & V & F & F & F & \textbf{V} \\
    V & F & V & V & V & \textbf{V} \\
    V & F & F & V & V & \textbf{V} \\
    F & V & V & V & V & \textbf{V} \\
    F & V & F & V & V & \textbf{V} \\
    F & F & V & V & V & \textbf{V} \\
    F & F & F & V & V & \textbf{V} \\
    \hline
    \end{tabular}
    \end{center}
    Es una \underline{\textbf{Tautología}}.

\end{enumerate}

\vspace{0.5cm}
\hrule
\vspace{0.5cm}